\documentclass[12pt]{article}
\usepackage[a3paper, margin=0.4in]{geometry}
\usepackage{amsmath}
\usepackage{xcolor}
\usepackage{amssymb}


\begin{document}

\section*{Domácí úkol IV}

Vypracoval: Daniel \textit{"Randál"} Ransdorf \hfill Podpis: \rule{4cm}{0.4pt}

\begin{enumerate}
  \item \textbf{Kuličky}

    Když vytáhneme 7 kuliček ze 30, očekáváme, že vytáhneme $\frac{1}{3}\cdot7 \approx 2.33$ skleněnek.
    Pravděpodobnosti vytažení skleněnek bez vracení budou vyšší kolem této očekávané hodnoty.
    Pravděpodobnosti na extrémech (blízko 0 a blízko 7) budou zase vyšší s vracením,
    protože při tahání bez vracení se postupně snižuje pravděpodobnost
    na takovýto "extrémní" výsledek.

    Případ vytažení tří skleněnek je bezpečně blízko očekávané hodnotě $\approx 2.33$,
    takže bude určitě vyšší pravděpodobnost vytáhnout tři skleněnky při tahání \textbf{bez vracení}.

  \item \textbf{Dvoustupňový screening}

    Označme jev A, že člověk užíval opium. Označme jev B, že výsledek vyjde pozitivní.
    Přepišme údaje ze zadání do pravděpodobností:
    \begin{align*}
      P(B \,|\, \neg A) &= 0.02 \quad \Rightarrow P(\neg B \,|\, \neg A) = 0.98 \\
      P(\neg B \,|\, A) &= 0.05 \quad \Rightarrow P(B \,|\, A) = 0.95 \\
      P(A) &= 0.01 \quad \Rightarrow P(\neg A) = 0.99
    \end{align*}
    Dopočtěme další hodnoty, které by mohly být užitečné
    \begin{align*}
      P(B \,|\, A) &= \frac{P(A \cap B)}{P(A)} = 0.95 \quad \Rightarrow P(A \cap B) = 0.95P(A) \\
    \end{align*}

    \begin{enumerate}
      \item Chceme zjistit $P(A \,|\, B^2)$
        \begin{align*}
          P(B^2 \,|\, A) &= P(B \,|\, A)^2 = 0.95^2 = 0.9025 \\
          P(B^2 \,|\, \neg A) &= P(B \,|\, \neg A)^2 = 0.02^2 = 0.0004 \\
          P(B^2) &= P(A) \cdot P(B^2 \,|\, A) + P(\neg A) \cdot P(B^2 \,|\, \neg A) = 0.01\cdot0.9025 + 0.99\cdot0.0004 = 0.009421 \\
          P(A \,|\, B^2) &= \frac{P(A)}{P(B^2)} \cdot P(B^2 \,|\, A) = \frac{0.01}{0.009421} \cdot 0.9025 \\
          P(A \,|\, B^2) &= \underline{\underline{0.957662456}}
        \end{align*}
      
      \item Vypočtěme $P(A \,|\, B)$ a poměry $\frac{P(A \,|\, B^2)}{P(A \,|\, B)}, \frac{P(B \,|\, \neg A)}{P(B^2 \,|\, \neg A)}, \frac{P(\neg B \,|\, A)}{P((\neg B)^2 \,|\, A)}$
        \begin{align*}
          P(B) &= P(A) \cdot P(B \,|\, A) + P(\neg A) \cdot P(B \,|\, \neg A) = 0.01\cdot0.95 + 0.99\cdot 0.02 \\
          P(B) &= 0.0293 \\
          P(A \,|\, B) &= \frac{P(A)}{P(B)} \cdot P(B \,|\, A) = \frac{0.01}{0.0293} \cdot 0.95 \\
          P(A \,|\, B) &= 0.3242320819 \\
          \frac{P(B \,|\, \neg A)}{P(B^2 \,|\, \neg A)} &= \frac{1}{P(B \,|\, \neg A)} = 0.02^{-1} = 50 \\
          \frac{P(\neg B \,|\, A)}{P((\neg B)^2 \,|\, A)} &= \frac{1}{P(\neg B \,|\, A)} = 0.05^{-1} = 20 \\
          \frac{P(A \,|\, B^2)}{P(A \,|\, B)} &= \frac{0.957662456}{0.3242320819} = 2.9536326276
        \end{align*}
      
      Použítím dvou testů místo jednoho se 20krát sníží podíl falesně negativních
      a 50krát sníží podíl falešně pozitivních. Když člověka pozitivně otestujeme
      podruhé, šance, že je uživatelem opia téměř 3krát vzroste, circa z 32\% na 96\%.
    \end{enumerate}

  \end{enumerate}



\end{document}
